\documentclass[11pt]{article}

\usepackage[utf8]{inputenc}
\usepackage[T1]{fontenc}
\usepackage{lmodern}
\usepackage{amsmath,amssymb}
\usepackage{graphicx}
\usepackage{booktabs}
\usepackage{array}
\usepackage{xcolor}
\usepackage[margin=1in]{geometry}
\usepackage{tikz}
\usetikzlibrary{arrows.meta,positioning,shapes.geometric}
\usepackage[hidelinks]{hyperref}
\usepackage{microtype}
\usepackage{enumitem}
\usepackage[round]{natbib}

\graphicspath{{../results/figures/}{./figures/}}

\newcommand{\zstar}{\ensuremath{Z^{\star}}}
\newcommand{\code}[1]{\texttt{#1}}
\newcommand{\Exec}{\ensuremath{\mathrm{Exec}}}

\title{\textbf{Sound Biosynthetic Design over Iterative Grammars:\\A Realizability Geometry for Fungal Polyketides}}
\author{Brage Eilertsen\\
\small \texttt{brageei@uio.no}\\
\small \textit{Affiliation to be completed}}
\date{May 2026 --- working draft (design reframe); \S6 first per the composition plan}

\begin{document}
\maketitle

% =====================================================================================
% DRAFTING ORDER (per the composition plan): S6 first (anchor), then S3 inheriting its
% phrasing, then S1 (intro), then S4/S5 (one page each), then S7, then the abstract last
% (the compression test for the whole paper). Sections other than S6 are stubs below.
% =====================================================================================

\begin{abstract}
% WRITTEN LAST -- the compression test. Seed: the framing-note compressed thesis. If the abstract
% writes easily, the spine held.
\textit{[Stub --- written last. Spine: wrong object ($\mathrm{BGC}\!\to\!$structure regression, the
51--68\% ceiling) $\to$ program space $y=\Exec_\Theta(z;G)$ run forward / inverse / backward $\to$ a
sound triage tool with named limits. Most of the abstract is the framing-note one-paragraph thesis.]}
\end{abstract}

\section{Introduction}\label{sec:intro}
\textit{[Stub --- drafted after \S6 and \S3. Wrong frame $\to$ right frame (program space; inference
\emph{and} design); the one-sentence thesis; architecture figure. \P2 lever: Cox (2023) identifies
machine learning as the route forward where rational HR-PKS engineering ``will remain extremely
difficult'' --- taken as a design \emph{constraint}, not a validation. \P3: structure elucidation here
proceeds by expression + NMR with MS as triage; we build a triage tool grounded in metabologenomics
observables and name explicitly where triage is insufficient.]}

\section{Substrate: a sound executor and verifier}\label{sec:substrate}
\textit{[Stub --- edited from the inverse-compiler manuscript (\code{paper/main.tex}). Sound
deterministic executor + beam verifier; the three-state verdict (\textsc{verified} /
\textsc{under-observed} / \textsc{out-of-grammar}); guarantees relative to operator library, grammar,
observables, mass tolerance, adduct model.]}

\section{Characterizing the natural manifold}\label{sec:manifold}
\textit{[Stub --- drafted second, inheriting \S\ref{sec:design}'s phrasing. Opener (does the reframe and
inoculates against being read as a limitations chapter): ``the natural fungal PKS manifold has 8
exactly-reachable cores, a 0.57 control-axis ceiling, and a tailoring-permissive interior with 64\%
conditional reconstructibility --- this section characterizes its shape.'' 8/326 = the manifold's
\emph{boundary condition} (not a prediction ceiling); the structural-similarity proxy critique as honest
defensive work ($8\!\to\!160$); the 0.57 wall as the axis along which the interesting designs live (not
where prediction fails). Refers forward to \S\ref{sec:design}.]}

\section{Generality across grammars}\label{sec:generality}
\textit{[Stub --- ONE PAGE. PKS, NRPS, and a typed PKS--NRPS hybrid behind one interface; the
load-bearing point is that the observability and verdict layers are family-agnostic, not the demo
counts. Reviewers wanting more cross-grammar work are pointed at the repository.]}

\section{Mining: genome and metabolome}\label{sec:mining}
\textit{[Stub --- ONE PAGE. The BGC$\leftrightarrow$peak assignment graph; structure verdict
$\perp$ attribution verdict; one real blind cluster$\to$mass case (6-MSA, MIBiG \code{BGC0001275}). A
bridge, not a pillar.]}

% =====================================================================================
% =====================================================================================
% Section 6 -- the design contribution. The anchor section: written to stand alone (as if S3 has not
% been drafted). Composed from the framing note's design paragraph, the realizability geometry, the
% literature curation, the worked 6-MSA figure, and the experiment-selection planner.
% =====================================================================================
\section{Sound biosynthetic design over the iteration grammar}\label{sec:design}

Running the executor forward turns a program into a structure; running the verifier inverse turns a
structure into the set $\zstar$ of programs that could have produced it. Here we run the grammar
\emph{backward}: given a desired, possibly non-natural target, we place it on a \emph{realizability
geometry} relative to the natural manifold, and we rank the experiments that would confirm it. Two
results make this more than a restatement of the diagnostic half. First, the design space decomposes
onto two axes that map to fundamentally different wet-lab capabilities. Second---and this is what keeps
the contribution honest---the cost on those axes is the field's \emph{measured} difficulty of fungal PKS
engineering, curated from primary literature, not an assumption we supply.

\subsection{The realizability geometry: two axes}\label{sec:realiz-geometry}
A design target is a program $z_d$; its product is $y_d=\Exec_\Theta(z_d;G)$. We score $z_d$ by its
edit distance to the nearest natural cluster on the manifold, decomposed into two typed axes:
\begin{itemize}[leftmargin=1.4em,itemsep=2pt]
  \item \textbf{Structural} edits change the \emph{domain content}: add a reductive domain
  (\code{add\_reductive}), swap the starter or extender (\code{starter}, \code{extender}), reprogram the
  release cyclase (\code{release}), or insert a methyltransferase de~novo (\code{add\_cmt}). These map
  onto bacterial modular-PKS engineering.
  \item \textbf{Control} edits change the \emph{iteration program} with the domain set fixed: which cycle
  a present domain fires on (\code{reduction}, \code{c\_methyl}) or the iteration count
  (\code{cycle\_add}, \code{cycle\_remove}). This is the iteration-grammar frontier---precisely the
  $100\%\!\to\!57\%$ wall of the cross-taxa transfer probe (\S\ref{sec:manifold}), now read as a design
  dimension rather than a prediction limit.
\end{itemize}
A design at coordinate $(2,1)$ (two structural, one control edit) is a different proposition from one at
$(3,0)$: the first requires reprogramming the iteration itself, the second is documented domain
engineering. Each target's most-realizable route is the minimum-cost edit path to a natural cluster, and
its \emph{verdict}---\textsc{natural}, \textsc{engineerable}, \textsc{frontier}, or
\textsc{speculative}---is read off that path. Crucially, the control axis \emph{only exists as a design
dimension} for a machine that distinguishes per-cycle firing within a fixed domain set; a tool without
the iteration grammar cannot represent the edit ``reprogram the ketoreductase at cycle~3,'' let alone
cost it (\S\ref{sec:moat}).

\subsection{A literature-grounded, super-additive cost}\label{sec:cost}
The design enumeration bounds the cost model to the nine edit kinds that actually appear in the 1-edit
neighbourhood, and we curated each from primary sources rather than scraping abstracts (two tier rulings
turned on facts a search snippet stated incorrectly). The structural edits inherit the mature
combinatorial-biosynthesis canon: acyltransferase/extender swaps and loading-module/starter swaps are
documented \citep{atreview,science1997}, as are reductive-loop transplants \citep{reductiveloop}; de~novo
C-methyltransferase insertion has no precedent (the methylation literature reprograms a \emph{present}
C-MeT, never inserts one) and is scored \textsc{speculative}.

The control edits are the unsolved frontier. Cox's review of iterative highly-reducing PKS programming
\citep{cox2023} states plainly that rational efforts ``will remain extremely difficult'' because the
programme is ``an emergent property\ldots highly unpredictable.'' Yet a \emph{single} edit beside a
closely related cluster is achievable: domain swaps between the tenellin and desmethylbassianin synthases
mapped methylation-pattern and chain-length changes cleanly \citep{fisch2011}. We encode exactly this
shape. Structural edits sum linearly by tier; control edits are \textbf{super-additive}---quadratic in
the number $k$ of control edits in a design:
\begin{equation}
\mathrm{cost}(z_d) \;=\; \underbrace{\textstyle\sum_{e\in\text{struct}} \tau(e)}_{\text{linear}}
\;+\; \underbrace{\Big(\textstyle\sum_{e\in\text{ctrl}} \tau(e)\Big)\cdot k_{\text{ctrl}}}_{\sim k^2,\ \text{super-additive}},
\label{eq:cost}
\end{equation}
where $\tau$ is the curated per-edit tier. Equation~\eqref{eq:cost} reduces to the linear tier cost at
$k=1$ (the Fisch--Cox single-edit regime) and grows quadratically as edits stack (the Cox emergent
regime). The cost model is therefore not asserted: it is the published difficulty landscape of fungal
PKS engineering, written into the objective. This is the difference between ``we built a design engine''
and ``we built a design engine whose cost model is the measured difficulty of the field.''

\subsection{The worked neighbourhood of 6-MSA}\label{sec:worked}
Figure~\ref{fig:worked} walks one clean template---the 6-methylsalicylic acid synthase, the cluster with
a verified blind real-mass match (\S\ref{sec:mining})---end to end. Each row of its 1-edit neighbourhood
is a build proposal carrying its $(\text{structural},\text{control})$ coordinate, realizability verdict,
predicted monoisotopic mass, and the observable that would confirm it. The neighbourhood separates
cleanly into an \textbf{engineerable} block (single structural edits---starter swaps such as
acetyl$\to$hexanoyl at mass~244.13, release reprograms---all confirmable by accurate mass) and a
\textbf{frontier} block (single control edits such as reprogramming the cycle-3 reduction), with
coverage-limited out-of-grammar designs sunk below the headline. One row reads, ``one documented edit
(acetyl$\to$hexanoyl) from the 6-MSA synthase, mass 244.13, verifiable by accurate mass''; another, ``one
iteration-program edit (reprogram cycle-3 reduction), mass 172.07, needs MS/MS.'' Design distance, the
predicted observable, and a verdict on whether triage suffices or structure elucidation is required---on
the same line.

\begin{figure}[t]
\centering
\fbox{\parbox{0.92\textwidth}{\small\itshape [Figure: the 1-edit design neighbourhood of 6-MSA, as a
ranked build table. Source: \code{scripts/worked\_design\_6msa.py}. Columns: $(S,C)$ edit decomposition;
verdict; predicted monoisotopic mass; verifying observable. Two labelled blocks (engineerable /
frontier), executor-verifiable designs leading.]}}
\caption{The worked 6-MSA design neighbourhood: design distance and verification plan on one line.}
\label{fig:worked}
\end{figure}

\subsection{Control-axis dominance}\label{sec:dominance}
Across the 1-edit neighbourhood of the natural manifold the control axis dominates: \textbf{69} designs
lie on the frontier (involve a control edit) against \textbf{49} engineerable (structural-only), with 24
speculative. That ratio is itself a product of the literature curation---the raw, uncurated count was a
starker $93\!:\!34$, and grounding the edit costs (demoting release to a structural domain swap, raising
de~novo C-MeT and iteration-count edits) softened it to a defensible margin. Fungal biosynthetic
\emph{design} is mostly iteration-program editing: the axis only a sound per-cycle executor can express.

\subsection{The experiment-selection planner}\label{sec:planner}
Given a design, which experiment should one run next? The planner ranks each modelled observation $o$ by
expected information gain per unit cost, weighted by realizability:
\begin{equation}
\mathrm{score}(o) \;=\; \mathbb{E}\!\left[\Delta|\zstar|\right]\cdot w_{\text{real}}\big/ c_{\text{obs}}(o),
\qquad w_{\text{real}} = 1\big/\mathrm{cost}(\text{unresolved edits}),
\label{eq:planner}
\end{equation}
with speculative designs hard-excluded ($w_{\text{real}}=0$). The two cost terms sit on different sides:
observation cost $c_{\text{obs}}$ (accurate mass $1$, MS/MS $3$, in committed relative units) in the
denominator, and the super-additive realizability cost of Eq.~\eqref{eq:cost} as the design-side weight.
Because that weight is super-additive, resolving one of two stacked control edits drops the residual
penalty non-linearly---so the planner reasons about the \emph{joint resolution state} of a design, not
merely the most-uncertain edge. The planner's output is a decision tree over observation outcomes in
which both children of every node are first-class: the refuted branch (the observation is inconsistent
with the predicted product) carries as much information as the confirmed one.

The planner returns a typed triage verdict. \textsc{Verified-by-triage}: the modelled observables
collapse $|\zstar|$ to near-unique. \textsc{Elucidation-required}: the modelled observables are
exhausted and $|\zstar|$ is still large---the design must go to structure elucidation.
\textsc{Out-of-grammar}: a mass observation refutes every legal program. \emph{Naming the
elucidation-required case is the contribution}: a triage tool that knows its own limits, and whose limits
are typed and machine-checkable. Figure~\ref{fig:trees} contrasts two trees: 6-MSA collapses at accurate
mass (a shallow tree---the planner doing its job), while the highly-reducing C$_{12}$ core BAB stays
under-determined through mass and MS/MS ($4050\!\to\!25\!\to\!23$ candidates) and is routed to expression
$+$ NMR (the planner naming its limit).

\begin{figure}[t]
\centering
\fbox{\parbox{0.92\textwidth}{\small\itshape [Figure: two decision trees side by side. Source:
\code{scripts/planner\_tree\_figure.py} $\to$ \code{results/planner\_tree\_\{6msa,bab\}.mmd}. Left:
6-MSA, root decision (run mass?) $\to$ confirmed $\Rightarrow$ verified-by-triage; refuted $\Rightarrow$
out-of-grammar. Right: BAB, mass $\to$ MS/MS exhausted $\Rightarrow$ elucidation-required; refuted leaves
drawn first-class at each level.]}}
\caption{The triage planner doing its job (6-MSA) and naming its own limit (BAB). Refuted branches are
rendered first-class, not collapsed.}
\label{fig:trees}
\end{figure}

This scoping is forced by the data, not chosen for modesty. All seven verified products in our corpus
had their genome$\to$structure link established by heterologous expression (with knockout or native
expression) and NMR, with accurate mass and MS/MS serving as triage and detection---never as the sole
structure-determining observable. The planner is therefore a metabologenomics \emph{triage} tool: it
answers ``which observation collapses $|\zstar|$ given the genome,'' not ``how is this structure
proven.'' Evaluated retrospectively against this curated corpus it routes correctly---the small aromatics
(orsellinic acid, 6-MSA, mellein, 6-hydroxymellein) verify by accurate mass, while the highly-reducing
core is flagged elucidation-required---and on the design neighbourhood it steers attention to
engineerable and frontier designs over speculative ones.

\subsection{The architectural moat and the same-wall symmetry}\label{sec:moat}
Two claims carry the contribution. \textbf{First, it is architectural, not a capability we happen to
have.} The control axis exists as a design dimension only for a sound executor that distinguishes
per-cycle firing within a fixed domain set. ``Reprogram the ketoreductase at cycle~3'' is not a hard
edit to a retrosynthesis tool (which works in chemical space), an enzyme-design tool (which works per
protein), or a pathway tool (which ignores assembly-line grammar)---it is a sentence about an object none
of them can represent. We are the missing layer between generative chemistry and protein design, and the
cost on the axis is the field's measured difficulty (Eq.~\eqref{eq:cost}). \textbf{Second, the same wall
has opposite valence in the two directions.} The $100\%\!\to\!57\%$ iteration-grammar collapse is, for
\emph{prediction}, a ceiling---where reconstruction fails (\S\ref{sec:manifold}); for \emph{design}, it
is the frontier---the dimension along which the interesting molecules live. One measurement, a limitation
read one way and the contribution read the other.

% =====================================================================================

\section{Limitations and outlook}\label{sec:limits}
\textit{[Stub --- drafted after \S6. v0 NRPS/hybrid grammars; literature-curated edit tiers (the
defensibility moat, curated from primary sources). \emph{Triage scoping, third placement:} the planner
selects experiments and names elucidation-required clusters --- it does not close the design--build--test
loop. Named v1 extensions: observation-updated \emph{probabilistic realizability} (tiers as success
distributions, needing wet-lab build-outcome data); the per-cluster mining planner; and
\emph{bacterial generalization} (ClusterCAD manifold density) as the post-paper companion piece, where
it lands as confirmation rather than competition.]}

\section{Conclusion}\label{sec:conclusion}
\textit{[Stub.]}

% ---- working bibliography (S6 refs; merges with main.tex's bibliography on assembly) ----
\begin{thebibliography}{99}\small
% Author-year labels match main.tex's natbib style. Authors for the structural-canon refs are to be
% completed from the verified PMIDs/DOIs in data/curation/edit_precedent.md -- NOT fabricated here.
\bibitem[Cox(2023)]{cox2023} R.~J. Cox. Curiouser and curiouser: progress in understanding the
programming of iterative highly-reducing polyketide synthases. \emph{Nat. Prod. Rep.}, 40:9--27, 2023.
doi:10.1039/D2NP00007E.
\bibitem[Fisch et al.(2011)]{fisch2011} K.~M. Fisch, W. Bakeer, A.~A. Yakasai, Z. Song, J. Pedrick,
Z. Wasil, A.~M. Bailey, C.~M. Lazarus, T.~J. Simpson, R.~J. Cox. Rational domain swaps decipher
programming in fungal highly-reducing polyketide synthases and resurrect an extinct metabolite.
\emph{J. Am. Chem. Soc.}, 133:16635--16641, 2011. doi:10.1021/ja206914q.
\bibitem[Zdouc et al.(2025)]{mibig} M.~M. Zdouc, K. Blin, et al. MIBiG 4.0: advancing biosynthetic gene
cluster curation. \emph{Nucleic Acids Res.}, 53:D678, 2025.
\bibitem[AT-engineering review(2018)]{atreview} Acyltransferases as tools for polyketide synthase
engineering. \emph{Antibiotics}, 7(3):62, 2018. PMC6164871. [authors to complete]
\bibitem[Erythromycin analogues(1997)]{science1997} Precursor-directed biosynthesis of erythromycin
analogs by an engineered polyketide synthase. \emph{Science}, 277:367, 1997.
doi:10.1126/science.277.5324.367. [authors to complete]
\bibitem[Reductive-loop swaps(2008)]{reductiveloop} A polylinker approach to reductive loop swaps in
modular polyketide synthases. PMID~18937219, 2008. [authors to complete]
\end{thebibliography}

\end{document}
